\subsubsection{Learning without Forgetting (LwF)}

\gls{lwf} is a distillation-based \gls{cl} method that mitigates \gls{cf}
without storing past-task samples. Instead of replay, it preserves previous
knowledge by matching the current model's outputs on old classes to those of a
frozen teacher model copied at task boundaries.

\paragraph{Knowledge Preservation by Distillation}

After completing task $t{-}1$, a snapshot $f_{\theta^*}$ of the model is frozen
as a teacher. During training on task $t$, the student model $f_\theta$ is
optimised with a standard supervised loss on current-task labels plus a
temperature-scaled distillation term on previously seen classes:
%
\begin{equation}\label{eqn:lwf_loss}
    \mathcal{L}_\text{LwF}
    =
    \mathcal{L}_\text{CE}^{(t)}
    +
    \lambda_\text{dist}
    \,T^2\,
    D_\text{KL}
    \!\left(
        \operatorname{softmax}\!\left(\frac{z_\text{teach}^{old}}{T}\right)
        \,\Big\|\,
        \operatorname{softmax}\!\left(\frac{z_\text{stud}^{old}}{T}\right)
    \right)
\end{equation}
%
where $z_\text{teach}^{old}$ and $z_\text{stud}^{old}$ are teacher and student
logits restricted to old classes, $T$ is the distillation temperature, and
$\lambda_\text{dist}$ controls the distillation strength.

\paragraph{\gls{cl} Approach}

In each optimisation step, \gls{lwf} performs two complementary objectives:
(i) fit the incoming task labels, and (ii) preserve old-task decision behavior
through logit matching to the frozen teacher. This transfers historical
knowledge through function-space constraints rather than through explicit sample
replay. The method therefore avoids replay memory overhead but depends on the
teacher outputs providing sufficient signal about previous tasks.

\paragraph{Inference}

\gls{lwf} uses the current model directly at inference time; no teacher model is
required. In \gls{til}, task identity is used to select the relevant logit
slice. In \gls{cil}, prediction is made from the full logit vector across all
seen classes.

